% Options for packages loaded elsewhere
\PassOptionsToPackage{unicode}{hyperref}
\PassOptionsToPackage{hyphens}{url}
\PassOptionsToPackage{dvipsnames,svgnames,x11names}{xcolor}
%
\documentclass[
]{article}
\usepackage{amsmath,amssymb}
\usepackage{lmodern}
\usepackage{iftex}
\ifPDFTeX
  \usepackage[T1]{fontenc}
  \usepackage[utf8]{inputenc}
  \usepackage{textcomp} % provide euro and other symbols
\else % if luatex or xetex
  \usepackage{unicode-math}
  \defaultfontfeatures{Scale=MatchLowercase}
  \defaultfontfeatures[\rmfamily]{Ligatures=TeX,Scale=1}
\fi
% Use upquote if available, for straight quotes in verbatim environments
\IfFileExists{upquote.sty}{\usepackage{upquote}}{}
\IfFileExists{microtype.sty}{% use microtype if available
  \usepackage[]{microtype}
  \UseMicrotypeSet[protrusion]{basicmath} % disable protrusion for tt fonts
}{}
\makeatletter
\@ifundefined{KOMAClassName}{% if non-KOMA class
  \IfFileExists{parskip.sty}{%
    \usepackage{parskip}
  }{% else
    \setlength{\parindent}{0pt}
    \setlength{\parskip}{6pt plus 2pt minus 1pt}}
}{% if KOMA class
  \KOMAoptions{parskip=half}}
\makeatother
\usepackage{xcolor}
\usepackage{color}
\usepackage{fancyvrb}
\newcommand{\VerbBar}{|}
\newcommand{\VERB}{\Verb[commandchars=\\\{\}]}
\DefineVerbatimEnvironment{Highlighting}{Verbatim}{commandchars=\\\{\}}
% Add ',fontsize=\small' for more characters per line
\newenvironment{Shaded}{}{}
\newcommand{\AlertTok}[1]{\textcolor[rgb]{1.00,0.00,0.00}{\textbf{#1}}}
\newcommand{\AnnotationTok}[1]{\textcolor[rgb]{0.38,0.63,0.69}{\textbf{\textit{#1}}}}
\newcommand{\AttributeTok}[1]{\textcolor[rgb]{0.49,0.56,0.16}{#1}}
\newcommand{\BaseNTok}[1]{\textcolor[rgb]{0.25,0.63,0.44}{#1}}
\newcommand{\BuiltInTok}[1]{\textcolor[rgb]{0.00,0.50,0.00}{#1}}
\newcommand{\CharTok}[1]{\textcolor[rgb]{0.25,0.44,0.63}{#1}}
\newcommand{\CommentTok}[1]{\textcolor[rgb]{0.38,0.63,0.69}{\textit{#1}}}
\newcommand{\CommentVarTok}[1]{\textcolor[rgb]{0.38,0.63,0.69}{\textbf{\textit{#1}}}}
\newcommand{\ConstantTok}[1]{\textcolor[rgb]{0.53,0.00,0.00}{#1}}
\newcommand{\ControlFlowTok}[1]{\textcolor[rgb]{0.00,0.44,0.13}{\textbf{#1}}}
\newcommand{\DataTypeTok}[1]{\textcolor[rgb]{0.56,0.13,0.00}{#1}}
\newcommand{\DecValTok}[1]{\textcolor[rgb]{0.25,0.63,0.44}{#1}}
\newcommand{\DocumentationTok}[1]{\textcolor[rgb]{0.73,0.13,0.13}{\textit{#1}}}
\newcommand{\ErrorTok}[1]{\textcolor[rgb]{1.00,0.00,0.00}{\textbf{#1}}}
\newcommand{\ExtensionTok}[1]{#1}
\newcommand{\FloatTok}[1]{\textcolor[rgb]{0.25,0.63,0.44}{#1}}
\newcommand{\FunctionTok}[1]{\textcolor[rgb]{0.02,0.16,0.49}{#1}}
\newcommand{\ImportTok}[1]{\textcolor[rgb]{0.00,0.50,0.00}{\textbf{#1}}}
\newcommand{\InformationTok}[1]{\textcolor[rgb]{0.38,0.63,0.69}{\textbf{\textit{#1}}}}
\newcommand{\KeywordTok}[1]{\textcolor[rgb]{0.00,0.44,0.13}{\textbf{#1}}}
\newcommand{\NormalTok}[1]{#1}
\newcommand{\OperatorTok}[1]{\textcolor[rgb]{0.40,0.40,0.40}{#1}}
\newcommand{\OtherTok}[1]{\textcolor[rgb]{0.00,0.44,0.13}{#1}}
\newcommand{\PreprocessorTok}[1]{\textcolor[rgb]{0.74,0.48,0.00}{#1}}
\newcommand{\RegionMarkerTok}[1]{#1}
\newcommand{\SpecialCharTok}[1]{\textcolor[rgb]{0.25,0.44,0.63}{#1}}
\newcommand{\SpecialStringTok}[1]{\textcolor[rgb]{0.73,0.40,0.53}{#1}}
\newcommand{\StringTok}[1]{\textcolor[rgb]{0.25,0.44,0.63}{#1}}
\newcommand{\VariableTok}[1]{\textcolor[rgb]{0.10,0.09,0.49}{#1}}
\newcommand{\VerbatimStringTok}[1]{\textcolor[rgb]{0.25,0.44,0.63}{#1}}
\newcommand{\WarningTok}[1]{\textcolor[rgb]{0.38,0.63,0.69}{\textbf{\textit{#1}}}}
\usepackage{graphicx}
\makeatletter
\def\maxwidth{\ifdim\Gin@nat@width>\linewidth\linewidth\else\Gin@nat@width\fi}
\def\maxheight{\ifdim\Gin@nat@height>\textheight\textheight\else\Gin@nat@height\fi}
\makeatother
% Scale images if necessary, so that they will not overflow the page
% margins by default, and it is still possible to overwrite the defaults
% using explicit options in \includegraphics[width, height, ...]{}
\setkeys{Gin}{width=\maxwidth,height=\maxheight,keepaspectratio}
% Set default figure placement to htbp
\makeatletter
\def\fps@figure{htbp}
\makeatother
\setlength{\emergencystretch}{3em} % prevent overfull lines
\providecommand{\tightlist}{%
  \setlength{\itemsep}{0pt}\setlength{\parskip}{0pt}}
\setcounter{secnumdepth}{-\maxdimen} % remove section numbering
% definitions for citeproc citations
\NewDocumentCommand\citeproctext{}{}
\NewDocumentCommand\citeproc{mm}{%
  \begingroup\def\citeproctext{#2}\cite{#1}\endgroup}
\makeatletter
 % allow citations to break across lines
 \let\@cite@ofmt\@firstofone
 % avoid brackets around text for \cite:
 \def\@biblabel#1{}
 \def\@cite#1#2{{#1\if@tempswa , #2\fi}}
\makeatother
\newlength{\cslhangindent}
\setlength{\cslhangindent}{1.5em}
\newlength{\csllabelwidth}
\setlength{\csllabelwidth}{3em}
\newenvironment{CSLReferences}[2] % #1 hanging-indent, #2 entry-spacing
 {\begin{list}{}{%
  \setlength{\itemindent}{0pt}
  \setlength{\leftmargin}{0pt}
  \setlength{\parsep}{0pt}
  % turn on hanging indent if param 1 is 1
  \ifodd #1
   \setlength{\leftmargin}{\cslhangindent}
   \setlength{\itemindent}{-1\cslhangindent}
  \fi
  % set entry spacing
  \setlength{\itemsep}{#2\baselineskip}}}
 {\end{list}}
\usepackage{calc}
\newcommand{\CSLBlock}[1]{\hfill\break\parbox[t]{\linewidth}{\strut\ignorespaces#1\strut}}
\newcommand{\CSLLeftMargin}[1]{\parbox[t]{\csllabelwidth}{\strut#1\strut}}
\newcommand{\CSLRightInline}[1]{\parbox[t]{\linewidth - \csllabelwidth}{\strut#1\strut}}
\newcommand{\CSLIndent}[1]{\hspace{\cslhangindent}#1}
\ifLuaTeX
\usepackage[bidi=basic]{babel}
\else
\usepackage[bidi=default]{babel}
\fi
\babelprovide[main,import]{american}
% get rid of language-specific shorthands (see #6817):
\let\LanguageShortHands\languageshorthands
\def\languageshorthands#1{}
\ifLuaTeX
  \usepackage{selnolig}  % disable illegal ligatures
\fi
\IfFileExists{bookmark.sty}{\usepackage{bookmark}}{\usepackage{hyperref}}
\IfFileExists{xurl.sty}{\usepackage{xurl}}{} % add URL line breaks if available
\urlstyle{same} % disable monospaced font for URLs
\hypersetup{
  pdftitle={HydraR: Stateful Agentic Orchestration for Scientific
Reproducibility in R},
  pdfauthor={Chi Nam Ignatius Pang, Aidan Tay},
  pdflang={en-US},
  colorlinks=true,
  linkcolor={Maroon},
  filecolor={Maroon},
  citecolor={Blue},
  urlcolor={Blue},
  pdfcreator={LaTeX via pandoc}}

\title{HydraR: Stateful Agentic Orchestration for Scientific
Reproducibility in R}

\definecolor{c53baa1}{RGB}{83,186,161}
\definecolor{c202826}{RGB}{32,40,38}
\def \rorglobalscale {0.1}
\newcommand{\rorlogo}{%
\begin{tikzpicture}[y=1cm, x=1cm, yscale=\rorglobalscale,xscale=\rorglobalscale, every node/.append style={scale=\rorglobalscale}, inner sep=0pt, outer sep=0pt]
  \begin{scope}[even odd rule,line join=round,miter limit=2.0,shift={(-0.025, 0.0216)}]
    \path[fill=c53baa1,nonzero rule,line join=round,miter limit=2.0] (1.8164, 3.012) -- (1.4954, 2.5204) -- (1.1742, 3.012) -- (1.8164, 3.012) -- cycle;
    \path[fill=c53baa1,nonzero rule,line join=round,miter limit=2.0] (3.1594, 3.012) -- (2.8385, 2.5204) -- (2.5172, 3.012) -- (3.1594, 3.012) -- cycle;
    \path[fill=c53baa1,nonzero rule,line join=round,miter limit=2.0] (1.1742, 0.0669) -- (1.4954, 0.5588) -- (1.8164, 0.0669) -- (1.1742, 0.0669) -- cycle;
    \path[fill=c53baa1,nonzero rule,line join=round,miter limit=2.0] (2.5172, 0.0669) -- (2.8385, 0.5588) -- (3.1594, 0.0669) -- (2.5172, 0.0669) -- cycle;
    \path[fill=c202826,nonzero rule,line join=round,miter limit=2.0] (3.8505, 1.4364).. controls (3.9643, 1.4576) and (4.0508, 1.5081) .. (4.1098, 1.5878).. controls (4.169, 1.6674) and (4.1984, 1.7642) .. (4.1984, 1.8777).. controls (4.1984, 1.9719) and (4.182, 2.0503) .. (4.1495, 2.1132).. controls (4.1169, 2.1762) and (4.0727, 2.2262) .. (4.0174, 2.2635).. controls (3.9621, 2.3006) and (3.8976, 2.3273) .. (3.824, 2.3432).. controls (3.7505, 2.359) and (3.6727, 2.367) .. (3.5909, 2.367) -- (2.9676, 2.367) -- (2.9676, 1.8688).. controls (2.9625, 1.8833) and (2.9572, 1.8976) .. (2.9514, 1.9119).. controls (2.9083, 2.0164) and (2.848, 2.1056) .. (2.7705, 2.1791).. controls (2.6929, 2.2527) and (2.6014, 2.3093) .. (2.495, 2.3487).. controls (2.3889, 2.3881) and (2.2728, 2.408) .. (2.1468, 2.408).. controls (2.0209, 2.408) and (1.905, 2.3881) .. (1.7986, 2.3487).. controls (1.6925, 2.3093) and (1.6007, 2.2527) .. (1.5232, 2.1791).. controls (1.4539, 2.1132) and (1.3983, 2.0346) .. (1.3565, 1.9436).. controls (1.3504, 2.009) and (1.3351, 2.0656) .. (1.3105, 2.1132).. controls (1.2779, 2.1762) and (1.2338, 2.2262) .. (1.1785, 2.2635).. controls (1.1232, 2.3006) and (1.0586, 2.3273) .. (0.985, 2.3432).. controls (0.9115, 2.359) and (0.8337, 2.367) .. (0.7519, 2.367) -- (0.1289, 2.367) -- (0.1289, 0.7562) -- (0.4837, 0.7562) -- (0.4837, 1.4002) -- (0.6588, 1.4002) -- (0.9956, 0.7562) -- (1.4211, 0.7562) -- (1.0118, 1.4364).. controls (1.1255, 1.4576) and (1.2121, 1.5081) .. (1.2711, 1.5878).. controls (1.2737, 1.5915) and (1.2761, 1.5954) .. (1.2787, 1.5991).. controls (1.2782, 1.5867) and (1.2779, 1.5743) .. (1.2779, 1.5616).. controls (1.2779, 1.4327) and (1.2996, 1.3158) .. (1.3428, 1.2113).. controls (1.3859, 1.1068) and (1.4462, 1.0176) .. (1.5237, 0.944).. controls (1.601, 0.8705) and (1.6928, 0.8139) .. (1.7992, 0.7744).. controls (1.9053, 0.735) and (2.0214, 0.7152) .. (2.1474, 0.7152).. controls (2.2733, 0.7152) and (2.3892, 0.735) .. (2.4956, 0.7744).. controls (2.6016, 0.8139) and (2.6935, 0.8705) .. (2.771, 0.944).. controls (2.8482, 1.0176) and (2.9086, 1.1068) .. (2.952, 1.2113).. controls (2.9578, 1.2253) and (2.9631, 1.2398) .. (2.9681, 1.2544) -- (2.9681, 0.7562) -- (3.3229, 0.7562) -- (3.3229, 1.4002) -- (3.4981, 1.4002) -- (3.8349, 0.7562) -- (4.2603, 0.7562) -- (3.8505, 1.4364) -- cycle(0.9628, 1.7777).. controls (0.9438, 1.7534) and (0.92, 1.7357) .. (0.8911, 1.7243).. controls (0.8623, 1.7129) and (0.83, 1.706) .. (0.7945, 1.7039).. controls (0.7588, 1.7015) and (0.7252, 1.7005) .. (0.6932, 1.7005) -- (0.4839, 1.7005) -- (0.4839, 2.0667) -- (0.716, 2.0667).. controls (0.7477, 2.0667) and (0.7805, 2.0643) .. (0.8139, 2.0598).. controls (0.8472, 2.0553) and (0.8768, 2.0466) .. (0.9025, 2.0336).. controls (0.9282, 2.0206) and (0.9496, 2.0021) .. (0.9663, 1.9778).. controls (0.9829, 1.9534) and (0.9914, 1.9209) .. (0.9914, 1.8799).. controls (0.9914, 1.8362) and (0.9819, 1.8021) .. (0.9628, 1.7777) -- cycle(2.6125, 1.3533).. controls (2.5889, 1.2904) and (2.5553, 1.2359) .. (2.5112, 1.1896).. controls (2.4672, 1.1433) and (2.4146, 1.1073) .. (2.3529, 1.0814).. controls (2.2916, 1.0554) and (2.2228, 1.0427) .. (2.1471, 1.0427).. controls (2.0712, 1.0427) and (2.0026, 1.0557) .. (1.9412, 1.0814).. controls (1.8799, 1.107) and (1.8272, 1.1433) .. (1.783, 1.1896).. controls (1.7391, 1.2359) and (1.7052, 1.2904) .. (1.6817, 1.3533).. controls (1.6581, 1.4163) and (1.6465, 1.4856) .. (1.6465, 1.5616).. controls (1.6465, 1.6359) and (1.6581, 1.705) .. (1.6817, 1.7687).. controls (1.7052, 1.8325) and (1.7388, 1.8873) .. (1.783, 1.9336).. controls (1.8269, 1.9799) and (1.8796, 2.0159) .. (1.9412, 2.0418).. controls (2.0026, 2.0675) and (2.0712, 2.0804) .. (2.1471, 2.0804).. controls (2.223, 2.0804) and (2.2916, 2.0675) .. (2.3529, 2.0418).. controls (2.4143, 2.0161) and (2.467, 1.9799) .. (2.5112, 1.9336).. controls (2.5551, 1.8873) and (2.5889, 1.8322) .. (2.6125, 1.7687).. controls (2.636, 1.705) and (2.6477, 1.6359) .. (2.6477, 1.5616).. controls (2.6477, 1.4856) and (2.636, 1.4163) .. (2.6125, 1.3533) -- cycle(3.8015, 1.7777).. controls (3.7825, 1.7534) and (3.7587, 1.7357) .. (3.7298, 1.7243).. controls (3.701, 1.7129) and (3.6687, 1.706) .. (3.6333, 1.7039).. controls (3.5975, 1.7015) and (3.5639, 1.7005) .. (3.5319, 1.7005) -- (3.3226, 1.7005) -- (3.3226, 2.0667) -- (3.5547, 2.0667).. controls (3.5864, 2.0667) and (3.6192, 2.0643) .. (3.6526, 2.0598).. controls (3.6859, 2.0553) and (3.7155, 2.0466) .. (3.7412, 2.0336).. controls (3.7669, 2.0206) and (3.7883, 2.0021) .. (3.805, 1.9778).. controls (3.8216, 1.9534) and (3.8301, 1.9209) .. (3.8301, 1.8799).. controls (3.8301, 1.8362) and (3.8206, 1.8021) .. (3.8015, 1.7777) -- cycle;
  \end{scope}
\end{tikzpicture}
}

%%%%%%%%%%%%%%%%%%%%%%%%%%%%%%%%%%%%%%%%%%%%%%%%%%%%%%%%%%%%%%%%%%%%%%%%
% Authors and Affiliations

\usepackage[affil-it]{authblk}
\usepackage{orcidlink}
%% \renewcommand\Authsep{, }
\setlength{\affilsep}{1em}
\author[1%
  %
  ]{Chi Nam Ignatius Pang%
    \,\orcidlink{0000-0001-9703-5741}\,%
    }
\author[1%
  %
  ]{Aidan Tay%
    \,\orcidlink{0000-0003-1315-4896}\,%
    }

\affil[1]{Australian Proteome Analysis Facility (APAF), Macquarie
University, Sydney, Australia%
  }
%%%%%%%%%%%%%%%%%%%%%%%%%%%%%%%%%%%%%%%%%%%%%%%%%%%%%%%%%%%%%%%%%%%%%%%%
\date{15 April 2026}

\begin{document}
\maketitle

\section{Summary}\label{summary}

\texttt{HydraR} is a lightweight, state-managed framework designed to
orchestrate complex ``agentic'' workflows---autonomous, multi-step
processes driven by Large Language Models (LLMs) natively within R.
While mainstream orchestration tools such as LangChain
(\citeproc{ref-langchain}{Chase, 2022}) and CrewAI
(\citeproc{ref-crewai}{Moura, 2023}) are predominantly Python-based,
\texttt{HydraR} fulfills the R community's critical requirements for
rigorous reproducibility, auditability, and seamless integration with
established statistical pipelines.

Researchers can architect workflows as Directed Acyclic Graphs (DAGs) or
iterative state machines, wherein each node represents an LLM-prompted
task, a deterministic R function, or an autonomous auditor.
\texttt{HydraR} ensures broad compatibility---supporting cloud APIs
(Gemini (\citeproc{ref-google2025gemini}{Google, 2025}), Claude, OpenAI)
and local models (Ollama)---while enforcing structural integrity through
a built-in \texttt{Advanced\ Validation\ Engine} that performs deep
static analysis of orchestration manifests before execution.

\section{Statement of Need}\label{statement-of-need}

The contemporary R ecosystem currently lacks a unified framework that
provides both high-level agentic orchestration and robust, low-level
filesystem safety. Existing packages, such as \texttt{ellmer}
(\citeproc{ref-ellmer}{Wickham et al., 2025}), primarily focus on
conversational interfaces, while \texttt{mall}
(\citeproc{ref-mall}{Couch, 2024}) excels at mapping LLM inferences over
textual data within data frames. Neither is structurally equipped to
manage long-running, multi-agent collaborations that demand complex
state persistence or isolated execution environments.

\texttt{HydraR} addresses three fundamental challenges in this domain:
1. \textbf{Auditability}: It implements a centralized
\texttt{AgentState} with persistent checkpointing (via DuckDB
(\citeproc{ref-duckdb}{Raasveldt \& Mühleisen, 2019})), yielding a
complete, verifiable audit trail of all LLM interactions and state
mutations. 2. \textbf{Filesystem Safety}: Leveraging Git Worktrees,
\texttt{HydraR} isolates parallel agent executions into temporary
branches, virtually eliminating data corruption and race conditions
during concurrent operations. 3. \textbf{Reproducible Orchestration}: By
defining workflows via \texttt{Mermaid.js} syntax and \texttt{YAML}
manifests, \texttt{HydraR} decouples orchestration logic from the
underlying execution framework. This declarative approach creates
portable, shareable agentic protocols, effortlessly bypassing the added
complexity typical when R is used to access multi-agent frameworks via
\texttt{reticulate} (\citeproc{ref-reticulate}{Ushey et al., 2024}).

\section{State of the Field}\label{state-of-the-field}

Although \texttt{reticulate} (\citeproc{ref-reticulate}{Ushey et al.,
2024}) enables R users to interface with Python frameworks like
LangChain, doing so introduces a structural impedance mismatch when
serializing complex R data objects across language boundaries.
\texttt{HydraR} avoids this by building natively on the \texttt{R6}
object-oriented system (\citeproc{ref-R6}{Chang, 2021}), preserving the
fidelity of statistical objects while allowing direct compatibility with
R's integrated debugging tool chains. Furthermore, unlike
\texttt{ellmer} (\citeproc{ref-ellmer}{Wickham et al., 2025}) (which
offers chat-centric workflows without DAG orchestration), and
\texttt{gptstudio} (\citeproc{ref-gptstudio}{Nivard \& Wade, 2024})
(which serves as an IDE-centric coding assistant), \texttt{HydraR} is
specifically engineered to manage the complete lifecycle of headless,
parallel, multi-agent collaborations replete with robust state
retention.

\section{Software Design}\label{software-design}

\texttt{HydraR} employs a highly modular \texttt{R6} architecture
spanning several key components: - \texttt{AgentLLMNode} /
\texttt{AgentLogicNode}: Encapsulate LLM-driven prompts and
deterministic pure-R logic steps, respectively. They utilize pluggable
\texttt{AgentDriver} backends for provider-agnostic model communication.
- \texttt{AgentDAG}: The core orchestration engine that manages node
dependencies, validates network topology (via \texttt{igraph}
(\citeproc{ref-igraph}{Csárdi \& Nepusz, 2006})), and executes logic in
parallel (via \texttt{furrr} (\citeproc{ref-furrr}{Vaughan \& Reznik,
2022})). It features an integrated compiler-like Validation Engine that
cross-references Mermaid topologies with \texttt{YAML} resource
definitions and lints embedded R logic for compliance with safety
standards (e.g., prohibiting imperative loops). - \texttt{AgentState}: A
centralized state repository employing versioned history and custom
reducers to systematically coordinate multi-agent outputs. -
\texttt{WorktreeManager}: An engine that provisions isolated Git
worktrees for parallel execution tasks, paired with an extensible
\texttt{ConflictResolver} handling automated semantic resolution
alongside human-in-the-loop task reconciliation.

\section{Research Applications}\label{research-applications}

\subsection{Multimodal Travel Itinerary
Planner}\label{multimodal-travel-itinerary-planner}

In this example, an LLM \texttt{Planner} agent generates a draft
itinerary, which is then audited by a deterministic \texttt{Validator}
against hard constraints. Successfully validated plans proceed to a
\emph{multimodal image generation} phase where \textbf{Gemini 3.1 Flash}
(\citeproc{ref-google2025gemini}{Google, 2025}) creates visual assets
based on the itinerary's locales. Finally, a logic node renders a
bespoke, CSS-styled HTML pamphlet. Crucially, \textbf{conditional
looping} between the planner and validator enforces an iterative
refinement cycle until all strict parameters are met
(\autoref{fig:travel_workflow}).

\protect\phantomsection\label{fig:travel_workflow}
\begin{figure}
\centering
\pandocbounded{\includegraphics[keepaspectratio]{figures/Itinerary Page 1.png}}
\caption{Page 1: Multi-modal travel pamphlet showcasing stateful
orchestration and image generation.}\label{fig:travel_workflow_a}
\end{figure}

\begin{figure}
\centering
\pandocbounded{\includegraphics[keepaspectratio]{figures/Itinerary Page 2.png}}
\caption{Page 2: Final formatted CSS-styled HTML
itinerary.}\label{fig:travel_workflow_b}
\end{figure}

Generated Multi-modal travel pamphlet for a Sydney to Hong Kong
itinerary.

\subsubsection{Declarative Workflow
Excerpt}\label{declarative-workflow-excerpt}

\begin{Shaded}
\begin{Highlighting}[]
\FunctionTok{graph}\KeywordTok{: }\CharTok{|}
\NormalTok{  graph LR}
\NormalTok{    Planner["Travel Planner | type=llm | role\_id=travel\_concierge"]}
\NormalTok{    Validator["Constraint Auditor | type=logic | logic\_id=validate\_constraints"]}
\NormalTok{    ImageGate["Image Gate | type=logic | logic\_id=check\_image\_status"]}
\NormalTok{    ImageGenerator["Image Generator | type=logic | logic\_id=generate\_and\_save\_images"]}
\NormalTok{    TemplateManager["Template Provider | type=logic | logic\_id=provide\_template"]}
\NormalTok{    PamphletFormatter["Pamphlet Formatter | type=logic | logic\_id=format\_pamphlet"]}
\NormalTok{    Finalizer["Itinerary Saver | type=logic | logic\_id=save\_itinerary"]}

\NormalTok{    Planner {-}{-}\textgreater{} Validator}
\NormalTok{    Validator {-}{-} "fail" {-}{-}\textgreater{} Planner}
\NormalTok{    Validator {-}{-} "pass" {-}{-}\textgreater{} ImageGate}
\NormalTok{    ImageGate {-}{-}\textgreater{} ImageGenerator}
\NormalTok{    ImageGate {-}{-}\textgreater{} TemplateManager}
\NormalTok{    ImageGenerator {-}{-}\textgreater{} TemplateManager}
\NormalTok{    TemplateManager {-}{-}\textgreater{} PamphletFormatter}
\NormalTok{    PamphletFormatter {-}{-}\textgreater{} Finalizer}

\FunctionTok{roles}\KeywordTok{:}
\FunctionTok{  travel\_concierge}\KeywordTok{: }\CharTok{\textgreater{}}
\NormalTok{    You are a professional travel concierge...}
\end{Highlighting}
\end{Shaded}

\begin{figure}
\centering
\pandocbounded{\includegraphics[keepaspectratio]{figures/travel_workflow.png}}
\caption{Visual DAG structure of the multi-modal travel itinerary
planner, featuring iterative validation loops and parallel generation
paths.}\label{fig:travel_graph}
\end{figure}

\subsubsection{R Orchestration (Minimal)}\label{r-orchestration-minimal}

\begin{Shaded}
\begin{Highlighting}[]
\FunctionTok{library}\NormalTok{(HydraR)}
\NormalTok{wf }\OtherTok{\textless{}{-}} \FunctionTok{load\_workflow}\NormalTok{(}\StringTok{"hong\_kong\_travel.yml"}\NormalTok{)}
\NormalTok{dag }\OtherTok{\textless{}{-}} \FunctionTok{spawn\_dag}\NormalTok{(wf, }\FunctionTok{auto\_node\_factory}\NormalTok{())}
\NormalTok{results }\OtherTok{\textless{}{-}}\NormalTok{ dag}\SpecialCharTok{$}\FunctionTok{run}\NormalTok{(}\AttributeTok{initial\_state =}\NormalTok{ wf}\SpecialCharTok{$}\NormalTok{initial\_state)}
\end{Highlighting}
\end{Shaded}

\subsection{Parallel Sorting Algorithm
Comparison}\label{parallel-sorting-algorithm-comparison}

This benchmarking example tasks three distinct LLM agents with
simultaneously implementing different sorting algorithms.
\texttt{HydraR} executes each task within an isolated
\texttt{Git\ worktree} to aggressively uncouple filesystem side-effects
and prevent merge conflicts. Following execution, a Merge Harmonizer
noded systematically reconciles the independent branches back to the
main state, preceding a terminal logic node that empirically benchmarks
the aggregated algorithms across five continuous trials
(\autoref{fig:sorting}).

\begin{figure}
\centering
\pandocbounded{\includegraphics[keepaspectratio]{figures/sorting_benchmark.pdf}}
\caption{Sorting Algorithm Performance Benchmark (1,000 elements over 5
trials).}\label{fig:sorting}
\end{figure}

\subsubsection{Declarative Workflow
Excerpt}\label{declarative-workflow-excerpt-1}

\begin{Shaded}
\begin{Highlighting}[]
\FunctionTok{graph}\KeywordTok{: }\CharTok{|}
\NormalTok{  graph LR}
\NormalTok{      bubble["Bubble Agent | type=llm | role\_id=bubble"]}
\NormalTok{      quick["Quick Agent | type=llm | role\_id=quick"]}
\NormalTok{      merge["Merge Agent | type=llm | role\_id=merge"]}
\NormalTok{      merger["Merge Harmonizer | type=merge"]}
\NormalTok{      benchmark["Benchmark | type=logic | logic\_id=run\_benchmark"]}

\NormalTok{      bubble {-}{-}\textgreater{} merger}
\NormalTok{      quick {-}{-}\textgreater{} merger}
\NormalTok{      merge {-}{-}\textgreater{} merger}
\NormalTok{      merger {-}{-}\textgreater{} benchmark}
\end{Highlighting}
\end{Shaded}

\begin{figure}
\centering
\pandocbounded{\includegraphics[keepaspectratio]{figures/sorting_workflow.png}}
\caption{Visual representation of the parallel sorting comparison
workflow, showing the fan-out/fan-in pattern and the merge
harmonizer.}\label{fig:sorting_graph}
\end{figure}

\subsubsection{R Orchestration (Parallel with
Isolation)}\label{r-orchestration-parallel-with-isolation}

\begin{Shaded}
\begin{Highlighting}[]
\FunctionTok{library}\NormalTok{(HydraR)}
\NormalTok{wf }\OtherTok{\textless{}{-}} \FunctionTok{load\_workflow}\NormalTok{(}\StringTok{"sorting\_benchmark.yml"}\NormalTok{)}
\NormalTok{dag }\OtherTok{\textless{}{-}} \FunctionTok{spawn\_dag}\NormalTok{(wf, }\FunctionTok{auto\_node\_factory}\NormalTok{())}
\NormalTok{results }\OtherTok{\textless{}{-}}\NormalTok{ dag}\SpecialCharTok{$}\FunctionTok{run}\NormalTok{(}\AttributeTok{initial\_state =}\NormalTok{ wf}\SpecialCharTok{$}\NormalTok{initial\_state, }\AttributeTok{use\_worktrees =} \ConstantTok{TRUE}\NormalTok{)}
\end{Highlighting}
\end{Shaded}

\subsection{Fault-Tolerant Pipelines with DuckDB
Persistence}\label{fault-tolerant-pipelines-with-duckdb-persistence}

To mitigate transient failures inherent to long-running scientific
workflows, \texttt{HydraR} integrates advanced checkpoint-and-resume
functionality powered by DuckDB (\citeproc{ref-duckdb}{Raasveldt \&
Mühleisen, 2019}). Whenever a pipeline pauses mid-execution, it
automatically persists its complete \texttt{AgentState}. Upon restart,
any operator-applied programmatic interventions are seamlessly merged,
actively preventing redundant evaluations by allowing execution to jump
directly back to the paused node.

\subsubsection{Declarative Workflow
Definition}\label{declarative-workflow-definition}

\begin{Shaded}
\begin{Highlighting}[]
\FunctionTok{graph}\KeywordTok{: }\CharTok{|}
\NormalTok{  graph LR}
\NormalTok{    Step1["Initialization"] {-}{-}\textgreater{} Step2["Risky Logic"]}
\NormalTok{    Step2["Risky Logic"] {-}{-}\textgreater{} Step3["Conclusion"]}
\FunctionTok{logic}\KeywordTok{:}
\FunctionTok{  check\_fixed}\KeywordTok{: }\CharTok{|}
\NormalTok{    if (!isTRUE(state$get("fixed"))) }
\NormalTok{      return(list(status = "PAUSE", output = "Waiting for manual fix."))}
\NormalTok{    list(status = "SUCCESS", output = "System recovered!")}
\end{Highlighting}
\end{Shaded}

\begin{figure}
\centering
\pandocbounded{\includegraphics[keepaspectratio]{figures/fault_workflow.png}}
\caption{Visual logic of the fault-tolerant pipeline with DuckDB
persistence, illustrating a linear flow with restartable
checkpoints.}\label{fig:fault_graph}
\end{figure}

\subsubsection{Checkpoint and Resume}\label{checkpoint-and-resume}

\begin{Shaded}
\begin{Highlighting}[]
\FunctionTok{library}\NormalTok{(HydraR)}
\NormalTok{wf    }\OtherTok{\textless{}{-}} \FunctionTok{load\_workflow}\NormalTok{(}\StringTok{"state\_persistence.yml"}\NormalTok{)}
\NormalTok{dag   }\OtherTok{\textless{}{-}} \FunctionTok{spawn\_dag}\NormalTok{(wf)}
\NormalTok{saver }\OtherTok{\textless{}{-}}\NormalTok{ DuckDBSaver}\SpecialCharTok{$}\FunctionTok{new}\NormalTok{(}\AttributeTok{db\_path =} \StringTok{"history.duckdb"}\NormalTok{)}
\NormalTok{tid   }\OtherTok{\textless{}{-}} \StringTok{"reprex{-}session{-}001"}

\CommentTok{\# Run 1: Pauses at Step2 — state saved to DuckDB}
\NormalTok{res1 }\OtherTok{\textless{}{-}}\NormalTok{ dag}\SpecialCharTok{$}\FunctionTok{run}\NormalTok{(}\AttributeTok{thread\_id =}\NormalTok{ tid, }\AttributeTok{checkpointer =}\NormalTok{ saver,}
                \AttributeTok{initial\_state =} \FunctionTok{list}\NormalTok{(}\AttributeTok{fixed =} \ConstantTok{FALSE}\NormalTok{))}

\CommentTok{\# Run 2: Fix applied, resume from Step2}
\NormalTok{res2 }\OtherTok{\textless{}{-}}\NormalTok{ dag}\SpecialCharTok{$}\FunctionTok{run}\NormalTok{(}\AttributeTok{thread\_id =}\NormalTok{ tid, }\AttributeTok{checkpointer =}\NormalTok{ saver,}
                \AttributeTok{initial\_state =} \FunctionTok{list}\NormalTok{(}\AttributeTok{fixed =} \ConstantTok{TRUE}\NormalTok{),}
                \AttributeTok{resume\_from =} \StringTok{"Step2"}\NormalTok{)}
\end{Highlighting}
\end{Shaded}

By safeguarding computational progress autonomously, this pattern
effectively circumvents prohibitively costly re-execution in data
science pipelines reliant on voluminous LLM invocations or laborious
large-scale database queries.

\section{Research Impact Statement}\label{research-impact-statement}

\texttt{HydraR} is actively deployed at the \textbf{Australian Proteome
Analysis Facility (APAF)} to architect and automate resilient
bioinformatics tooling. Its declarative YAML schemas empower
non-programming domain experts to routinely govern advanced LLM logic
pipelines without any R coding proficiency, thereby democratizing the
scale and reproducibility of model-assisted scientific research.

\section{AI Usage Disclosure}\label{ai-usage-disclosure}

Transparency and accountability are core tenets of the \texttt{HydraR}
paradigm. A substantial portion of the implementation was co-developed
utilizing \textbf{Antigravity}, an agentic AI coding assistant,
supervised by a rigorous ``human-in-the-loop'' review protocol.
Specifically, the original authors independently conceptualized the
underlying system architecture, while the AI generated specialized
logic, unit tests, and software documentation with every resultant
commit sequentially tested and manually verified.

\section{Acknowledgements}\label{acknowledgements}

We acknowledge funding from \textbf{Bioplatforms Australia}, enabled by
the \textbf{National Collaborative Research Infrastructure Strategy
(NCRIS)}.

\section*{References}\label{references}
\addcontentsline{toc}{section}{References}

\protect\phantomsection\label{refs}
\begin{CSLReferences}{1}{0}
\bibitem[\citeproctext]{ref-R6}
Chang, W. (2021). \emph{R6: Recursive object-oriented class system}.
\url{https://R6.r-lib.org}

\bibitem[\citeproctext]{ref-langchain}
Chase, H. (2022). LangChain: Building applications with LLMs through
composability. \emph{GitHub Repository}.
\url{https://github.com/langchain-ai/langchain}

\bibitem[\citeproctext]{ref-mall}
Couch, S. P. (2024). \emph{Mall: Run language models multiple times}.
\url{https://github.com/simonpcouch/mall}

\bibitem[\citeproctext]{ref-igraph}
Csárdi, G., \& Nepusz, T. (2006). The igraph software package for
complex network research. \emph{InterJournal}, \emph{Complex Systems},
1695. \url{https://igraph.org}

\bibitem[\citeproctext]{ref-google2025gemini}
Google. (2025). \emph{Gemini 1.5 flash}.
\url{https://gemini.google.com}.

\bibitem[\citeproctext]{ref-crewai}
Moura, J. (2023). CrewAI: Collaborative AI agents. \emph{GitHub
Repository}. \url{https://github.com/joaomoura/crewai}

\bibitem[\citeproctext]{ref-gptstudio}
Nivard, M., \& Wade, J. (2024). \emph{Gptstudio: Use large language
models directly in your development environment}.
\url{https://github.com/MichelNivard/gptstudio}

\bibitem[\citeproctext]{ref-duckdb}
Raasveldt, M., \& Mühleisen, H. (2019). DuckDB: An embeddable analytical
database. \emph{Proceedings of the 2019 International Conference on
Management of Data (SIGMOD)}.
\url{https://doi.org/10.1145/3299869.3320212}

\bibitem[\citeproctext]{ref-reticulate}
Ushey, K., Allaire, J., \& Tang, Y. (2024). \emph{Reticulate: Interface
to 'python'}. \url{https://rstudio.github.io/reticulate/}

\bibitem[\citeproctext]{ref-furrr}
Vaughan, D., \& Reznik, D. (2022). \emph{Furrr: Apply mapping functions
in parallel using futures}. \url{https://github.com/DavisVaughan/furrr}

\bibitem[\citeproctext]{ref-ellmer}
Wickham, H., Cheng, J., Jacobs, A., Aden-Buie, G., \& Schloerke, B.
(2025). \emph{Ellmer: Chat with large language models}.
\url{https://ellmer.tidyverse.org}

\end{CSLReferences}

\end{document}
